\documentclass[aspectratio=169,10pt]{beamer}

% Shared Beamer setup for Fall 2026 course slides.
% Clean Cal Poly-inspired theme: no custom class files or uncommon packages.
\mode<presentation>{
  \usetheme{default}
  \usefonttheme{professionalfonts}
}

\definecolor{CalPolyGreen}{HTML}{154734}
\definecolor{CalPolyGold}{HTML}{BD8B13}
\definecolor{CalPolySage}{HTML}{7FA28F}
\definecolor{CalPolyMint}{HTML}{DDF1E7}
\definecolor{CalPolySand}{HTML}{A29F86}
\definecolor{CalPolyBeige}{HTML}{EFEEDE}
\definecolor{DeckInk}{HTML}{1F2933}
\definecolor{DeckMuted}{HTML}{5F6B7A}
\definecolor{DeckSoft}{HTML}{F7F7F5}

\setbeamersize{text margin left=0.55in,text margin right=0.55in}
\setbeamercolor{normal text}{fg=DeckInk,bg=white}
\setbeamercolor{structure}{fg=CalPolyGreen}
\setbeamercolor{title}{fg=white}
\setbeamercolor{frametitle}{fg=CalPolyGreen,bg=white}
\setbeamercolor{block title}{bg=CalPolySage,fg=white}
\setbeamercolor{block body}{bg=CalPolyMint,fg=black}
\setbeamercolor{alerted text}{fg=CalPolyGold}

\setbeamerfont{title}{size=\Huge,series=\bfseries}
\setbeamerfont{frametitle}{size=\Large,series=\bfseries}
\setbeamerfont{block title}{size=\normalsize,series=\bfseries}

\setbeamertemplate{navigation symbols}{}
\setbeamertemplate{itemize item}{\color{CalPolyGold}$\blacktriangleright$}
\setbeamertemplate{itemize subitem}{\color{CalPolyGreen}$\bullet$}
\setbeamertemplate{footline}{
  \hfill{\color{DeckMuted}\scriptsize\insertframenumber/\inserttotalframenumber}
  \hspace{0.18in}\vspace{0.12in}
}
\setbeamertemplate{frametitle}{
  \vspace{0.18in}
  \hspace{0.02in}\insertframetitle\par
  \vspace{0.08in}
  \noindent\makebox[\linewidth][l]{%
    {\color{CalPolyGold}\rule{0.55in}{2pt}}%
    {\color{CalPolyGreen}\rule{\dimexpr\linewidth-0.55in\relax}{2pt}}%
  }
}

\newcommand{\CourseNumber}{Course}
\newcommand{\CourseTitle}{Course Title}
\newcommand{\LectureNumber}{0}
\newcommand{\LectureTitle}{Lecture Title}
\newcommand{\LectureDate}{Date}
\newcommand{\CourseUnit}{Unit}

\newcommand{\coursenumber}[1]{\renewcommand{\CourseNumber}{#1}}
\newcommand{\coursetitle}[1]{\renewcommand{\CourseTitle}{#1}}
\newcommand{\lecturenumber}[1]{\renewcommand{\LectureNumber}{#1}}
\newcommand{\lecturetitle}[1]{\renewcommand{\LectureTitle}{#1}}
\newcommand{\lecturedate}[1]{\renewcommand{\LectureDate}{#1}}
\newcommand{\courseunit}[1]{\renewcommand{\CourseUnit}{#1}}

\author{Kirk Duran}
\institute{California Polytechnic State University, San Luis Obispo}
\date{\LectureDate}

\newcommand{\makedecktitle}{
  \begingroup
  \setbeamercolor{background canvas}{bg=CalPolyGreen}
  \begin{frame}[plain]
    \vfill
    \begin{center}
      {\usebeamerfont{title}\color{white}\LectureTitle\par}
      \vspace{0.18in}
      {\Large\color{white}Lecture \LectureNumber\par}
    \end{center}
    \vfill
    \noindent{\color{white}\rule{\linewidth}{1.2pt}}\par
    \vspace{0.08in}
    \noindent\colorbox{CalPolyGold}{\makebox[0.24\linewidth][c]{\color{white}\Large\CourseNumber}}
    \hspace{0.12in}{\color{white}\Large Kirk Duran}
  \end{frame}
  \endgroup
}

\newenvironment{keyidea}{\begin{block}{Key idea}}{\end{block}}
\newenvironment{activity}{\begin{block}{In class}}{\end{block}}
\newenvironment{cpdefinition}{\begin{block}{Definition}}{\end{block}}
\newenvironment{exampleblockcp}{
  \begingroup
  \setbeamercolor{block title}{bg=CalPolySand,fg=white}
  \setbeamercolor{block body}{bg=CalPolyBeige,fg=black}
  \begin{block}{Example}
}{
  \end{block}
  \endgroup
}

\newcommand{\imageplaceholder}[2][0.3in]{%
  \noindent\makebox[\linewidth][c]{%
    \fcolorbox{CalPolySage}{CalPolyBeige}{%
      \parbox[c][#1][c]{0.86\linewidth}{%
        \centering
        {\color{DeckMuted}\scriptsize Image placeholder: }%
        {\color{DeckInk}\scriptsize\ttfamily\detokenize{#2}}%
      }%
    }%
  }%
}

\newcommand{\sectionframe}[1]{
  \begingroup
  \setbeamercolor{background canvas}{bg=CalPolyGreen}
  \begin{frame}[plain]
    \vfill
    {\color{white}\LARGE\bfseries #1\par}
    \vspace{0.18in}
    {\color{CalPolyGold}\rule{0.22\paperwidth}{3pt}}
    {\color{white}\rule{0.62\paperwidth}{3pt}}
    \vfill
  \end{frame}
  \endgroup
}

\newcommand{\placeholdernote}{
  \begin{block}{Placeholder}
    Replace these bullets with the final lecture material, examples, and in-class activities.
  \end{block}
}


\pdfstringdefDisableCommands{\def\translate#1{#1}}

\graphicspath{{../assets/lecture-22/}}

\newcommand{\lectureimage}[2][1.75in]{%
  \includegraphics[width=\linewidth,height=#1,keepaspectratio]{#2}%
}

\setbeamersize{text margin left=0.42in,text margin right=0.42in}

\coursenumber{CS 1001}
\coursetitle{Introduction to Computing}
\lecturenumber{22}
\lecturetitle{Iteration I --- \texttt{for} and \texttt{range}}
\lecturedate{Fall 2026}
\courseunit{9. Iteration}

% Python semantics checked against:
% https://docs.python.org/3/reference/compound_stmts.html#the-for-statement
% https://docs.python.org/3/library/stdtypes.html#ranges
\begin{document}

\makedecktitle

\begin{frame}{Today}
  \begin{itemize}
    \item Explain why repeated computation needs a new control-flow structure.
    \item Trace a \texttt{for} loop one iteration at a time.
    \item Treat the loop variable as a name that is repeatedly rebound.
    \item Use \texttt{range} to describe integer sequences for iteration.
    \item Distinguish direct iteration over values from iteration over indexes.
    \item Introduce an \textbf{accumulator} as state updated across iterations.
  \end{itemize}

  \begin{keyidea}
    A loop does not execute ``all at once.'' It repeatedly binds a loop variable, executes the body, and then advances to the next value.
  \end{keyidea}
\end{frame}

\sectionframe{Part 1: Why Programs Need Iteration}

\begin{frame}{Repeated Statements Do Not Scale}
  Suppose we want to display four values from a list.

  \begin{block}{Without iteration}
    \texttt{print(values[0])}\\
    \texttt{print(values[1])}\\
    \texttt{print(values[2])}\\
    \texttt{print(values[3])}
  \end{block}

  \begin{itemize}
    \item The same operation is written repeatedly.
    \item The code is tied to one specific list length.
    \item A larger input requires more source code.
  \end{itemize}

  \begin{keyidea}
    Iteration lets the \emph{program} repeat a computation instead of requiring the programmer to repeat source code.
  \end{keyidea}
\end{frame}

\begin{frame}{Definition: Iteration}
  \textbf{Iteration} is repeated execution of a block of code according to a control structure.
  \begin{columns}[c,onlytextwidth]
    \begin{column}{0.46\textwidth}
      For today's \texttt{for} loops:
      \begin{enumerate}
        \item Get the next supplied value, if one exists.
        \item Bind the loop variable to it.
        \item Execute the body, then request the next value.
      \end{enumerate}
      When no values remain, continue after the loop.
    \end{column}
    \begin{column}{0.49\textwidth}
      \lectureimage[1.95in]{for-loop-cycle.png}
    \end{column}
  \end{columns}
  Expressions, assignments, calls, and conditionals retain their usual meanings.
\end{frame}

\begin{frame}{Anatomy of a \texttt{for} Loop}
  \begin{block}{General form}
    \texttt{for variable in sequence:}\\
    \hspace{0.20in}\texttt{statement}\\
    \hspace{0.20in}\texttt{statement}
  \end{block}

  \begin{center}
    \renewcommand{\arraystretch}{1.25}
    \begin{tabular}{l|l}
      \textbf{Part} & \textbf{Role} \\\hline
      \texttt{for} & begins the loop \\
      loop variable & receives one value at a time \\
      sequence & supplies values to the loop \\
      indented body & executes once for each supplied value
    \end{tabular}
  \end{center}
\end{frame}

\begin{frame}{The Execution Model}
  Consider:

  \begin{block}{Code}
    \texttt{for value in [4, 9, 2]:}\\
    \hspace{0.20in}\texttt{print(value)}
  \end{block}

  Python conceptually performs:
  \begin{enumerate}
    \item bind \texttt{value} to \texttt{4}; execute the body;
    \item bind \texttt{value} to \texttt{9}; execute the body;
    \item bind \texttt{value} to \texttt{2}; execute the body;
    \item no values remain, so execution continues after the loop.
  \end{enumerate}

  \begin{keyidea}
    One loop body can execute many times, but each individual execution still follows ordinary statement-by-statement evaluation.
  \end{keyidea}
\end{frame}

\begin{frame}{Trace a Simple Loop}
  \begin{block}{Code}
    \texttt{values = [4, 9, 2]}\\
    \texttt{for value in values:}\\
    \hspace{0.20in}\texttt{answer = value + 1}\\
    \hspace{0.20in}\texttt{print(answer)}
  \end{block}

  \begin{activity}
    Before revealing the trace, predict the three values printed and the sequence of bindings for \texttt{value} and \texttt{answer}.
  \end{activity}
\end{frame}

\begin{frame}{Trace Table}
  \begin{center}
    \renewcommand{\arraystretch}{1.35}
    \begin{tabular}{c|c|c|c}
      \textbf{Iteration} & \textbf{\texttt{value}} & \textbf{\texttt{value + 1}} & \textbf{\texttt{answer}} \\\hline
      1 & 4 & 5 & 5 \\
      2 & 9 & 10 & 10 \\
      3 & 2 & 3 & 3
    \end{tabular}
  \end{center}

  \begin{block}{Printed output}
    \texttt{5}\\
    \texttt{10}\\
    \texttt{3}
  \end{block}

  The body is evaluated three separate times using three different bindings for \texttt{value}.
\end{frame}

\begin{frame}{The Loop Variable Is a Normal Binding}
  \begin{columns}[c,onlytextwidth]
    \begin{column}{0.43\textwidth}
      \begin{block}{Code}
        \texttt{for value in [4, 9, 2]:}\\
        \hspace{0.18in}\texttt{print(value)}
      \end{block}
      The picture shows three moments in time. Each supplied element takes over the \texttt{value} binding; the list stays unchanged.
    \end{column}
    \begin{column}{0.52\textwidth}
      \lectureimage[1.8in]{loop-variable-rebinding.png}
    \end{column}
  \end{columns}
  \begin{keyidea}
    Each iteration \textbf{rebinds} the loop variable to the next supplied value. It does not replace a position in the list.
  \end{keyidea}
\end{frame}

\begin{frame}{Indentation Defines the Loop Body}
  \begin{columns}[T,onlytextwidth]
    \begin{column}{0.47\textwidth}
      \begin{block}{Inside the loop}
        \texttt{for x in [1, 2, 3]:}\\
        \hspace{0.18in}\texttt{print(x)}\\
        \hspace{0.18in}\texttt{print("again")}
      \end{block}

      Both prints execute three times.
    \end{column}
    \begin{column}{0.47\textwidth}
      \begin{block}{After the loop}
        \texttt{for x in [1, 2, 3]:}\\
        \hspace{0.18in}\texttt{print(x)}\\
        \texttt{print("done")}
      \end{block}

      \texttt{"done"} prints once.
    \end{column}
  \end{columns}

  \begin{alertblock}{Tracing rule}
    Before tracing a loop, mark exactly which statements belong to its indented body.
  \end{alertblock}
\end{frame}

\sectionframe{Part 2: Iterating Directly Over Sequential Data}

\begin{frame}{Direct Iteration Over Values}
  Lists already contain the values we often want to process.

  \begin{block}{Code}
    \texttt{temperatures = [68, 71, 74]}\\
    \texttt{for temperature in temperatures:}\\
    \hspace{0.20in}\texttt{print(temperature)}
  \end{block}

  The loop variable receives the list elements directly:

  \begin{center}
    \texttt{68} $\rightarrow$ \texttt{71} $\rightarrow$ \texttt{74}
  \end{center}

  \begin{keyidea}
    If a computation needs the values themselves, direct iteration usually expresses that intent most clearly.
  \end{keyidea}
\end{frame}

\begin{frame}{Prediction: Expressions Inside a Loop}
  \begin{block}{Code}
    \texttt{values = [2, 5, 8]}\\
    \texttt{for value in values:}\\
    \hspace{0.20in}\texttt{result = value * 2}\\
    \hspace{0.20in}\texttt{print(result)}\\
    \texttt{print("done")}
  \end{block}

  \begin{activity}
    Predict the complete output. How many times is \texttt{value * 2} evaluated? How many times is \texttt{"done"} printed?
  \end{activity}
\end{frame}

\begin{frame}{Prediction: Reveal}
  \begin{block}{Output}
    \texttt{4}\\
    \texttt{10}\\
    \texttt{16}\\
    \texttt{done}
  \end{block}

  \begin{itemize}
    \item The loop body executes once for each of the three list elements.
    \item \texttt{result} is rebound during each iteration.
    \item The final print is outside the loop, so it executes once.
  \end{itemize}
\end{frame}

\begin{frame}{Rebinding the Loop Variable Does Not Mutate the List}
  \begin{block}{Code}
    \texttt{values = [2, 5, 8]}\\
    \texttt{for value in values:}\\
    \hspace{0.20in}\texttt{value = value * 2}\\
    \texttt{print(values)}
  \end{block}
  \textbf{Output:} \texttt{[2, 5, 8]}
  Each multiplication rebinds \texttt{value} to a new integer. The next iteration then binds it to the next original element.
  \begin{keyidea}
    As in Lecture 21, rebinding a name differs from mutating a list. This code never assigns to \texttt{values[index]}.
  \end{keyidea}
\end{frame}

\begin{frame}{The Same Mechanism Applies to Other Sequences}
  A \texttt{for} loop can receive values from sequential data other than lists.

  \begin{block}{Example}
    \texttt{for character in "cat":}\\
    \hspace{0.20in}\texttt{print(character)}
  \end{block}

  The bindings are:
  \begin{center}
    \texttt{character} $\rightarrow$ \texttt{"c"}\qquad
    \texttt{"a"}\qquad
    \texttt{"t"}
  \end{center}

  \begin{alertblock}{For now}
    Today the emphasis is the loop execution model. We will use iteration over lists and strings more systematically in the next lecture.
  \end{alertblock}
\end{frame}

\begin{frame}{A Loop Inside a Function}
  \begin{block}{Code}
    \texttt{def display(values: list[int]) -> None:}\\
    \hspace{0.20in}\texttt{for value in values:}\\
    \hspace{0.40in}\texttt{print(value)}
  \end{block}

  Function-call reasoning still applies:
  \begin{enumerate}
    \item evaluate the argument;
    \item bind the parameter \texttt{values};
    \item execute the body: the loop repeatedly binds \texttt{value} and prints it;
    \item the call ends after the loop completes.
  \end{enumerate}

  \begin{keyidea}
    Iteration happens \emph{inside} the existing call frame; it does not create a new function call for each iteration.
  \end{keyidea}
\end{frame}

\sectionframe{Part 3: Integer Sequences with \texttt{range}}

\begin{frame}{Why \texttt{range}?}
  Sometimes we want integer values that are not already stored in a list.

  \begin{block}{Example goal}
    Repeat a computation for the integers \texttt{0}, \texttt{1}, \texttt{2}, and \texttt{3}.
  \end{block}

  Python's \texttt{range} describes a sequence of integers suitable for iteration.

  \begin{block}{Code}
    \texttt{for i in range(4):}\\
    \hspace{0.20in}\texttt{print(i)}
  \end{block}
\end{frame}

\begin{frame}{\texttt{range(stop)}}
  \begin{block}{General form}
    \texttt{range(stop)}
  \end{block}

  It begins at \texttt{0} and stops \textbf{before} \texttt{stop}.

  \begin{center}
    \renewcommand{\arraystretch}{1.25}
    \begin{tabular}{c|l}
      \textbf{Expression} & \textbf{Integer values supplied} \\\hline
      \texttt{range(4)} & \texttt{0, 1, 2, 3} \\
      \texttt{range(1)} & \texttt{0} \\
      \texttt{range(6)} & \texttt{0, 1, 2, 3, 4, 5}
    \end{tabular}
  \end{center}
\end{frame}

\begin{frame}{The Stop Value Is Exclusive}
  You have already seen this convention in slicing.

  \begin{center}
    \texttt{values[1:4]} \qquad indexes \texttt{1, 2, 3}\\[0.18in]
    \texttt{range(1, 4)} \qquad integers \texttt{1, 2, 3}
  \end{center}

  \begin{keyidea}
    In both slicing and \texttt{range}, the start is included and the stop is excluded.
  \end{keyidea}

  This makes \texttt{range(len(values))} naturally line up with valid list indexes from \texttt{0} through \texttt{len(values) - 1}.
\end{frame}

\begin{frame}{\texttt{range(start, stop)}}
  \begin{block}{General form}
    \texttt{range(start, stop)}
  \end{block}

  \begin{center}
    \renewcommand{\arraystretch}{1.25}
    \begin{tabular}{c|l}
      \textbf{Expression} & \textbf{Integer values supplied} \\\hline
      \texttt{range(2, 6)} & \texttt{2, 3, 4, 5} \\
      \texttt{range(5, 8)} & \texttt{5, 6, 7} \\
      \texttt{range(3, 4)} & \texttt{3}
    \end{tabular}
  \end{center}

  The default step is \texttt{1}.
\end{frame}

\begin{frame}{\texttt{range(start, stop, step)}}
  \begin{block}{General form}
    \texttt{range(start, stop, step)}
  \end{block}

  \begin{center}
    \renewcommand{\arraystretch}{1.25}
    \begin{tabular}{c|l}
      \textbf{Expression} & \textbf{Integer values supplied} \\\hline
      \texttt{range(0, 8, 2)} & \texttt{0, 2, 4, 6} \\
      \texttt{range(1, 10, 3)} & \texttt{1, 4, 7} \\
      \texttt{range(2, 9, 2)} & \texttt{2, 4, 6, 8}
    \end{tabular}
  \end{center}

  \begin{keyidea}
    The step controls how the current integer changes before the next iteration.
  \end{keyidea}
\end{frame}

\begin{frame}{Negative Steps}
  A negative step describes decreasing integer values.

  \begin{center}
    \renewcommand{\arraystretch}{1.25}
    \begin{tabular}{c|l}
      \textbf{Expression} & \textbf{Integer values supplied} \\\hline
      \texttt{range(4, 0, -1)} & \texttt{4, 3, 2, 1} \\
      \texttt{range(6, 1, -2)} & \texttt{6, 4, 2}
    \end{tabular}
  \end{center}

  The stop remains exclusive even when the sequence decreases.
\end{frame}


\begin{frame}{A Valid \texttt{range} Can Be Empty}
  \begin{center}
    \renewcommand{\arraystretch}{1.35}
    \begin{tabular}{l|l}
      \textbf{Expression} & \textbf{Why it supplies no values} \\\hline
      \texttt{range(0)} & Start equals stop. \\
      \texttt{range(3, 3)} & Start equals stop. \\
      \texttt{range(5, 2)} & Positive step, but start is already past stop. \\
      \texttt{range(1, 5, -1)} & Negative step moves away from stop.
    \end{tabular}
  \end{center}
  An empty sequence means the body executes \textbf{zero times}; execution continues after the loop.
  \begin{alertblock}{Different case: an invalid step}
    \texttt{range(0, 5, 0)} raises \texttt{ValueError}. A zero step is invalid, not an empty range.
  \end{alertblock}
\end{frame}

\begin{frame}{Predict the Values from \texttt{range}}
  Without running Python, write the values supplied by each expression:

  \begin{enumerate}
    \item \texttt{range(5)}
    \item \texttt{range(2, 7)}
    \item \texttt{range(1, 8, 2)}
    \item \texttt{range(5, 1, -1)}
  \end{enumerate}

  \begin{activity}
    For each one, identify the start, stop, and step before listing the values.
  \end{activity}
\end{frame}


\begin{frame}{\texttt{range} Practice: Reveal}
  \begin{center}
    \renewcommand{\arraystretch}{1.5}
    \begin{tabular}{l|r|r|r|l}
      \textbf{Expression} & \textbf{Start} & \textbf{Stop} & \textbf{Step} & \textbf{Values supplied} \\\hline
      \texttt{range(5)} & 0 & 5 & 1 & \texttt{0, 1, 2, 3, 4} \\
      \texttt{range(2, 7)} & 2 & 7 & 1 & \texttt{2, 3, 4, 5, 6} \\
      \texttt{range(1, 8, 2)} & 1 & 8 & 2 & \texttt{1, 3, 5, 7} \\
      \texttt{range(5, 1, -1)} & 5 & 1 & $-1$ & \texttt{5, 4, 3, 2}
    \end{tabular}
  \end{center}
  \begin{keyidea}
    Identify the defaults first. Stop remains exclusive for both increasing and decreasing progressions.
  \end{keyidea}
\end{frame}

\begin{frame}{Trace a Loop Driven by \texttt{range}}
  \begin{block}{Code}
    \texttt{for i in range(1, 5):}\\
    \hspace{0.20in}\texttt{value = i * 10}\\
    \hspace{0.20in}\texttt{print(value)}
  \end{block}

  \begin{center}
    \renewcommand{\arraystretch}{1.3}
    \begin{tabular}{c|c|c}
      \textbf{Iteration} & \textbf{\texttt{i}} & \textbf{\texttt{value}} \\\hline
      1 & 1 & 10 \\
      2 & 2 & 20 \\
      3 & 3 & 30 \\
      4 & 4 & 40
    \end{tabular}
  \end{center}
\end{frame}

\begin{frame}{A \texttt{range} Describes the Integers}
  \begin{columns}[c,onlytextwidth]
    \begin{column}{0.42\textwidth}
      \begin{block}{Code}
        \texttt{numbers = range(0, 8, 2)}
      \end{block}
      It supplies \texttt{0, 2, 4, 6}.
      \begin{itemize}
        \item Start at \texttt{0}.
        \item Advance by \texttt{2}.
        \item Stop before \texttt{8}.
      \end{itemize}
    \end{column}
    \begin{column}{0.53\textwidth}
      \lectureimage[2.0in]{range-exclusive-stop.png}
    \end{column}
  \end{columns}
  \begin{keyidea}
    A \texttt{range} describes an integer progression without constructing a list of all its values.
  \end{keyidea}
\end{frame}

\begin{frame}{Using \texttt{range(len(...))} for Indexes}
  Sometimes the computation needs an element's \textbf{index}.

  \begin{block}{Code}
    \texttt{values = [8, 3, 6]}\\
    \texttt{for idx in range(len(values)):}\\
    \hspace{0.20in}\texttt{print(idx, values[idx])}
  \end{block}

  Since \texttt{len(values)} is \texttt{3}:
  \begin{center}
    \texttt{range(3)} $\Rightarrow$ \texttt{0, 1, 2}
  \end{center}

  Those are exactly the valid indexes for the list.
\end{frame}

\begin{frame}{Values or Indexes?}
  \begin{columns}[T,onlytextwidth]
    \begin{column}{0.47\textwidth}
      \begin{block}{Need each value}
        \texttt{for value in values:}\\
        \hspace{0.18in}\texttt{print(value)}
      \end{block}

      Loop variable: \textbf{element value}
    \end{column}
    \begin{column}{0.47\textwidth}
      \begin{block}{Need each index}
        \texttt{for idx in range(len(values)):}\\
        \hspace{0.18in}\texttt{print(idx)}
      \end{block}

      Loop variable: \textbf{integer index}
    \end{column}
  \end{columns}

  \begin{keyidea}
    Do not introduce indexes unless the computation actually needs the position. Direct iteration is simpler when only the values matter.
  \end{keyidea}
\end{frame}

\sectionframe{Part 4: State That Changes Across Iterations}

\begin{frame}{Some Computations Need Memory Across Iterations}
  \begin{columns}[c,onlytextwidth]
    \begin{column}{0.43\textwidth}
      Each iteration adds one value from \texttt{[4, 9, 2, 1]} to the sum so far.
      \begin{block}{Required state}
        \texttt{total}: the sum of the values processed so far.
      \end{block}
    \end{column}
    \begin{column}{0.52\textwidth}
      \lectureimage[1.8in]{accumulator-running-total.png}
    \end{column}
  \end{columns}
  \begin{keyidea}
    The boxes show successive values of \textbf{one} variable, not separate accumulators. Each update uses the previous total.
  \end{keyidea}
\end{frame}

\begin{frame}{Definition: Accumulator}
  \begin{cpdefinition}
    An \textbf{accumulator} is a variable that stores a partial result and is updated repeatedly as a loop executes.
  \end{cpdefinition}

  A common structure is:

  \begin{block}{Pattern}
    initialize accumulator\\
    \texttt{for value in sequence:}\\
    \hspace{0.20in}update accumulator using \texttt{value}\\
    use the completed accumulator after the loop
  \end{block}

  Today we use summation only to understand the mechanics. We will study additional loop patterns next lecture.
\end{frame}

\begin{frame}{The Lab Example: \texttt{tally}}
  \begin{block}{Code}
    \texttt{def tally(nums: list[int]) -> int:}\\
    \hspace{0.20in}\texttt{total = 0}\\
    \hspace{0.20in}\texttt{for num in nums:}\\
    \hspace{0.40in}\texttt{total = total + num}\\
    \hspace{0.20in}\texttt{return total}
  \end{block}

  For the call:
  \begin{center}
    \texttt{tally([4, 9, 2, 1])}
  \end{center}

  \begin{activity}
    Predict the values bound to \texttt{num} and the value of \texttt{total} after each iteration.
  \end{activity}
\end{frame}

\begin{frame}{Tracing the Accumulator}
  \begin{center}
    \renewcommand{\arraystretch}{1.35}
    \begin{tabular}{c|c|c|c}
      \textbf{Iteration} & \textbf{\texttt{num}} & \textbf{\texttt{total} before} & \textbf{\texttt{total} after} \\\hline
      -- & -- & -- & 0 \\
      1 & 4 & 0 & 4 \\
      2 & 9 & 4 & 13 \\
      3 & 2 & 13 & 15 \\
      4 & 1 & 15 & 16
    \end{tabular}
  \end{center}

  \begin{keyidea}
    The next iteration begins with the accumulator value produced by the previous iteration.
  \end{keyidea}
\end{frame}

\begin{frame}{Initialization Determines the Starting State}
  Why does \texttt{tally} begin with:
  \begin{center}
    \texttt{total = 0}?
  \end{center}

  Before any elements have been processed, the sum-so-far should be zero.

  \begin{itemize}
    \item initialization happens once, before the loop;
    \item updating happens once per iteration;
    \item the final value is used after all iterations complete.
  \end{itemize}

  \begin{alertblock}{Tracing discipline}
    Record the initial value before tracing. With an empty input, \texttt{tally([])} returns \texttt{0}: no updates occur.
  \end{alertblock}
\end{frame}

\begin{frame}{Return Placement Matters}
  Both fragments are inside a function, after \texttt{total = 0}.
  \begin{columns}[T,onlytextwidth]
    \begin{column}{0.48\textwidth}
      \begin{block}{After the loop}
        \texttt{for num in nums:}\\
        \hspace{0.18in}\texttt{total = total + num}\\
        \texttt{return total}
      \end{block}

      All elements are processed before returning.
    \end{column}
    \begin{column}{0.48\textwidth}
      \begin{block}{Inside the loop}
        \texttt{for num in nums:}\\
        \hspace{0.18in}\texttt{total = total + num}\\
        \hspace{0.18in}\texttt{return total}
      \end{block}

      For nonempty input, the first iteration returns immediately. Empty input reaches the function end and returns \texttt{None}.
    \end{column}
  \end{columns}

  \begin{keyidea}
    \texttt{return} keeps its old meaning: it immediately ends the current function call, even when it appears inside a loop.
  \end{keyidea}
\end{frame}

\sectionframe{Part 5: Integrating the Model}

\begin{frame}{Prediction Challenge}
  \begin{block}{Code}
    \texttt{def compute(start: int, stop: int) -> int:}\\
    \hspace{0.20in}\texttt{total = 0}\\
    \hspace{0.20in}\texttt{for value in range(start, stop):}\\
    \hspace{0.40in}\texttt{total = total + value}\\
    \hspace{0.20in}\texttt{return total}\\[0.08in]
    \texttt{answer = compute(2, 6)}
  \end{block}

  \begin{activity}
    Determine the values successively bound to \texttt{value}, the value of \texttt{total} after each iteration, and the final value bound to \texttt{answer}.
  \end{activity}
\end{frame}


\begin{frame}{Prediction Challenge: Reveal}
  For \texttt{compute(2, 6)}, the loop receives \texttt{2, 3, 4, 5}.
  \begin{center}
    \renewcommand{\arraystretch}{1.35}
    \begin{tabular}{c|c|c|c}
      \textbf{Iteration} & \textbf{\texttt{value}} & \textbf{\texttt{total} before} & \textbf{\texttt{total} after} \\\hline
      1 & 2 & 0 & 2 \\
      2 & 3 & 2 & 5 \\
      3 & 4 & 5 & 9 \\
      4 & 5 & 9 & 14
    \end{tabular}
  \end{center}
  \begin{block}{After the loop}
    \texttt{return total} returns \texttt{14}, so \texttt{answer} is bound to \texttt{14}.
  \end{block}
  \begin{activity}
    What would \texttt{compute(6, 2)} return? Explain using the range and the initial total.
  \end{activity}
  % Answer: 0, because range(6, 2) is empty with its default positive step.
\end{frame}

\begin{frame}{Where We Stop Today}
  Today's responsibility is the mechanics of iteration:
  \begin{itemize}
    \item how a \texttt{for} loop repeatedly executes;
    \item how the loop variable is rebound;
    \item how \texttt{range} supplies integer values;
    \item when to iterate over values versus indexes;
    \item how one accumulator can preserve state across iterations.
  \end{itemize}

  \begin{alertblock}{Next lectures}
    \textbf{Lecture 23:} common loop patterns such as counting, building results, and combining loops with conditionals.\\[0.06in]
    \textbf{Lecture 24:} higher-level sequence patterns, including comprehensions and map/filter-style reasoning.
  \end{alertblock}
\end{frame}

\begin{frame}{Takeaways}
  \begin{enumerate}
    \item A \texttt{for} loop processes one supplied value at a time.
    \item The loop variable is repeatedly rebound to the next value.
    \item The indented body executes once per supplied value.
    \item \texttt{range} describes integer values using start, exclusive stop, and optional step.
    \item Use direct iteration when you need values; use indexes when you need positions.
    \item An accumulator stores a partial result across iterations.
    \item Existing rules still apply: expressions evaluate normally, assignment rebinds names, and \texttt{return} immediately ends a function call.
  \end{enumerate}
\end{frame}

\end{document}
